% Correlated analytic representation of alpha: action, convergence, and
% intervals.
\subsubsection{Analytic representation of the alpha coordinate}
\label{subsec:alpha-lector-completo-rev08}

The control of the preceding jet requires the tail to be exhibited: merely
naming a complete reader is not sufficient. Let
\(\mathsf s\in\mathcal X_\alpha^{\mathrm{enr}}\). Before decimal
normalization, the same state already expresses the precoordinate
\begin{equation}
 a(\mathsf s):=p(\mathsf s)+e_0(\mathsf s)-f(\mathsf s)
                 -\kappa_{\mathrm{per}}(\mathsf s)\in I_\alpha,
 \label{eq:alpha-precoordenada-comun}
\end{equation}
with \(p=\pi_{\mathrm{HMT}}-3\), \(e_0=e_{\mathrm{HMT}}-2\), and
\(f=\varphi_{\mathrm{HMT}}-1\). Equivalently,
\(a=\pi_{\mathrm{HMT}}+e_{\mathrm{HMT}}-
\varphi_{\mathrm{HMT}}-4-\kappa_{\mathrm{per}}\). No conventional
\(\alpha\) is introduced here: \(p,e_0,f\) are the three regional
characters already generated, and \(\kappa_{\mathrm{per}}\) is the rational
readout of the terminal record. The integral closure applies to
\eqref{eq:alpha-precoordenada-comun} the normalization of
\eqref{eq:alpha-recurrencia}; the chart that follows expresses a
compatible analytic representation on the same precoordinate. Its dependency
graph receives \((p,e_0,f,\kappa_{\mathrm{per}})\) directly: it contains no
edge \(\alpha_A\to\lambda\), although the equality proved below permits the
two representation names to be identified subsequently.

Write \(q=1/729\), let \(F_{\mathsf s}\) be the ordered field generated by
the coordinates of \(\mathsf s\), and set
\(E_{9,\mathsf s}(X)=E_{21/22}^{(9)}(\mathsf s;X)\). Define the linking
coefficient
\begin{equation}
 \lambda(\mathsf s):=
 -\frac{E_{9,\mathsf s}(a(\mathsf s))
         \bigl(1-q\,a(\mathsf s)^3\bigr)}{a(\mathsf s)^{10}}.
 \label{eq:alpha-lambda-estado}
\end{equation}
All its arguments belong to the state before metrological recognition. In
particular, \(\lambda\) is not obtained by fitting a target decimal string.
Equation~\eqref{eq:alpha-lambda-estado} must be read with its exact type: once
the first representation of the state has produced \(a(\mathsf s)\), it fixes
the unique prolongation in the geometric class defined below. It does not
constitute a second independent selection of \(a(\mathsf s)\), but rather a
correlated analytic representation of the same enriched state.

Indeed, for \(\mu\in F_{\mathsf s}\), set
\[
 E_{\mathsf s,\mu}(X)
 :=E_{9,\mathsf s}(X)+\mu\frac{X^{10}}{1-qX^3}.
\]
Since \(0<a(\mathsf s)<10^{-2}\) and \(qa(\mathsf s)^3<1\), one has
\begin{equation}
 E_{\mathsf s,\mu}(a(\mathsf s))=0
 \quad\Longleftrightarrow\quad
 \mu=\lambda(\mathsf s).
 \label{eq:alpha-unicidad-lambda-clase-geometrica}
\end{equation}
Thus, the coefficient seed contains no freedom once the state and the class
of \(q\)-geometric prolongations have been fixed. This uniqueness is relative
to the declared class; it does not assert that the ninth-order jet by itself
selects a tail among all possible analytic series.

The local action on blocks of three coefficients is
\begin{equation}
 \mathcal C_\alpha:F_{\mathsf s}^{3}\longrightarrow F_{\mathsf s}^{3},
 \qquad \mathcal C_\alpha(u,v,w)=q(u,v,w),
 \qquad c^{(0)}=(\lambda(\mathsf s),0,0).
 \label{eq:alpha-accion-coeficientes}
\end{equation}
Therefore, for every \(n\geq0\),
\begin{equation}
 b_{10+3n}=q^n\lambda(\mathsf s),\qquad
 b_{11+3n}=b_{12+3n}=0.
 \label{eq:alpha-recurrencia-coeficientes}
\end{equation}
This is the rule that computes every coefficient beyond degree nine. No free
parameter remains in each prolongation.

For \(M\geq0\), let
\begin{equation}
 E_{\mathsf s,M}(X):=E_{9,\mathsf s}(X)
 +\lambda(\mathsf s)\sum_{n=0}^{M}q^nX^{10+3n}.
 \label{eq:alpha-truncamientos-completos}
\end{equation}
The complete function is
\begin{equation}
 E_{\mathsf s,\infty}(X)
 :=E_{9,\mathsf s}(X)
 +\lambda(\mathsf s)\frac{X^{10}}{1-X^3/729}.
 \label{eq:alpha-funcion-completa-explicita}
\end{equation}

To formulate the projective system without implicit types, fix
\(d_M=12+3M\) and define the jet fibration
\begin{equation}
 \mathfrak J_M:=
 \coprod_{\mathsf s\in\mathcal X_\alpha^{\mathrm{enr}}}
 \{\mathsf s\}\times F_{\mathsf s}[X]/(X^{d_M+1}).
 \label{eq:alpha-espacio-jets-rev08}
\end{equation}
If \(M'\ge M\), the restriction is
\begin{equation}
 \tau_{M',M}(\mathsf s,[P]_{d_{M'}})
 :=(\mathsf s,[P]_{d_M}).
 \label{eq:alpha-truncamiento-jets-rev08}
\end{equation}

\begin{proposition}[Naturality of the analytic recurrence]
\label{prop:alpha-naturalidad-recurrencia-explicita}
Let
\[
 \Gamma_M(\mathsf s)
 =\bigl(\mathsf s,[E_{\mathsf s,M}]_{d_M}\bigr).
\]
If \(M'\geq M\), the truncation \(\tau_{M',M}\) satisfies
\begin{equation}
 \tau_{M',M}\circ\Gamma_{M'}=\Gamma_M,
 \qquad
 \Gamma_{E21}(\mathsf s)=\varprojlim_M\Gamma_M(\mathsf s).
 \label{eq:alpha-naturalidad-explicita}
\end{equation}
The projection onto the three new degrees is precisely
\(\mathcal C_\alpha^{M+1}(c^{(0)})\).
\end{proposition}

\begin{proof}
Equation~\eqref{eq:alpha-recurrencia-coeficientes} fixes the block of degrees
\(10+3n,11+3n,12+3n\). Passing from \(M\) to \(M+1\) adds only the next
block; truncating it recovers the preceding polynomial literally.
Compatibility for \(M'\geq M\) follows by composition. Thus, the inverse
limit is not a section chosen from the bare jet, but the unique orbit of
\(c^{(0)}\) under \(\mathcal C_\alpha\).
\end{proof}

\subsubsection{Uniform convergence, simple root, and rational intervals}
\label{subsec:alpha-cotas-completas-rev08}

The regional cylinders and the periodic record provide bounds through
rational interval arithmetic. For each coordinate
\[
 c\in\{\pi_{\mathrm{HMT}},e_{\mathrm{HMT}},\varphi_{\mathrm{HMT}}\},
\]
let \(P_c\) be its certified thousand-digit nonadic prefix. The datum used is
not the rational number \(P_c\) as though it were the complete value, but the
half-open cylinder and its outer closure
\[
 \mathcal C_c^{(1000)}=[P_c,P_c+\varepsilon),
 \qquad
 \overline{\mathcal C}_c^{(1000)}=[P_c,P_c+\varepsilon],
 \qquad \varepsilon=10^{-1000}.
\]
Directed arithmetic operates conservatively with these closures. If
\(P_\pi,P_e,P_\varphi\) denote the three prefixes, one obtains
\begin{equation}
\begin{aligned}
 a_-&:=P_\pi+P_e-(P_\varphi+\varepsilon)-4-
       \kappa_{\mathrm{per}},\\
 a_+&:=(P_\pi+\varepsilon)+(P_e+\varepsilon)-P_\varphi-4-
       \kappa_{\mathrm{per}},\\
 A_{1000}&:=[a_-,a_+],
 \qquad a(\mathsf s)\in[a_-,a_+]=A_{1000},
 \qquad \operatorname{diam}A_{1000}=3\varepsilon.
\end{aligned}
\label{eq:alpha-propagacion-colas-rev08}
\end{equation}
The natural interval extension of \eqref{eq:alpha-lambda-estado} then
transports \(A_{1000}\), the cylinder of \(\pi_{\mathrm{HMT}}\), and the
directed rational bound on \(\delta\) to a closed interval
\(\Lambda_{1000}\ni\lambda(\mathsf s)\); none of the three tails is frozen.
In particular,
\begin{equation}
\begin{aligned}
 \frac{7297352569283800997285105472380662}{10^{36}}
 &<a(\mathsf s)<
 \frac{7297352569283800997285105472380663}{10^{36}},\\
 -8.711\mathbin{\times}10^{-6}
 &<\lambda(\mathsf s)<-8.709\mathbin{\times}10^{-6}.
\end{aligned}
 \label{eq:alpha-cotas-precoordenada-lambda}
\end{equation}
These inequalities use certified prefixes with their rational tail bound;
they do not receive an alpha table. The package verifier reproduces them from
the three nonadic outputs and from the terminal chart.

Let
\[
 r:=\frac{10^{-6}}{729}<1.372\mathbin{\times}10^{-9}.
\]
For \(0\leq x\leq10^{-2}\), the tail beyond \(M\) satisfies
\begin{equation}
 \left|E_{\mathsf s,\infty}(x)-E_{\mathsf s,M}(x)\right|
 \leq
 8.711\mathbin{\times}10^{-26}\,\frac{r^{M+1}}{1-r}.
 \label{eq:alpha-cota-cola-uniforme}
\end{equation}
Moreover,
\begin{equation}
 \left|\frac{\mathrm d}{\mathrm dx}
  \left(\lambda\frac{x^{10}}{1-x^3/729}\right)\right|
 \leq 8.711\mathbin{\times}10^{-6}
 \left(\frac{10\,10^{-18}}{1-r}
 +\frac{3\,10^{-24}/729}{(1-r)^2}\right)
 <8.712\mathbin{\times}10^{-23}.
 \label{eq:alpha-cota-derivada-cola}
\end{equation}
The tail of the derivatives beyond block \(M\) also admits the
depth-dependent estimate
\begin{equation}
 \sup_{0\le x\le10^{-2}}
 \left|E'_{\mathsf s,\infty}(x)-E'_{\mathsf s,M}(x)\right|
 \le 8.711\mathbin{\times}10^{-24}\,r^{M+1}
 \left(\frac{13+3M}{1-r}+\frac{3r}{(1-r)^2}\right)
 \xrightarrow[M\to\infty]{}0.
 \label{eq:alpha-cota-resto-derivadas-rev08}
\end{equation}

\begin{theorem}[Limit function and unique simple root of the correlated chart]
\label{thm:alpha-raiz-completa-explicita}
The sequence \(E_{\mathsf s,M}\), together with its derivatives, converges
uniformly to \(E_{\mathsf s,\infty}\) on
\(\overline I_\alpha=[0,10^{-2}]\). One has
\begin{equation}
 E_{\mathsf s,\infty}(a(\mathsf s))=0,
 \qquad
 \inf_{x\in\overline I_\alpha}
 E'_{\mathsf s,\infty}(x)>6.239.
 \label{eq:alpha-raiz-y-derivada-completas}
\end{equation}
Consequently, \(a(\mathsf s)\) is the unique simple root of the complete
function in \(I_\alpha\).
\end{theorem}

\begin{proof}
The ratio of two consecutive terms of the tail is bounded by \(r<1\), which
gives \eqref{eq:alpha-cota-cola-uniforme}; term-by-term differentiation and
the sum of the two geometric series give
\eqref{eq:alpha-cota-derivada-cola} and
\eqref{eq:alpha-cota-resto-derivadas-rev08}. The proof for the jet establishes
\(E'_{9,\mathsf s}>6.24\) on \(\overline I_\alpha\); therefore, the complete
perturbation leaves the derivative above \(6.239\). Finally, substituting
\eqref{eq:alpha-lambda-estado} into
\eqref{eq:alpha-funcion-completa-explicita}, the two summands evaluated at
\(a(\mathsf s)\) cancel exactly. The function is strictly increasing; hence
its unique root is simple and coincides with that precoordinate.
\end{proof}

To make the identification effective at arbitrary depth, directed evaluations
at the endpoints first certify
\[
 E_{\mathsf s,\infty}(0)<0<
 E_{\mathsf s,\infty}(10^{-2}),
 \qquad
 E'_{\mathsf s,\infty}\geq\mu:=6239/1000>c:=6
 \quad\hbox{on }[0,10^{-2}].
\]
By continuity there is a root in the interval and, by the second inequality,
it is unique. Only after this existence, monotonicity, and the invariant that
the bracket contains the root have been fixed is the mean-value-theorem
rescue applied. Start from \(D_0=[0,10^{-2}]\). If
\(D_j=[r_j,s_j]\), let \(m_j=(r_j+s_j)/2\), and let
\([F_j^-,F_j^+]\) be a directed rational enclosure of
\(E_{\mathsf s,\infty}(m_j)\), obtained by propagating the cylinders of
\(\pi_{\mathrm{HMT}},e_{\mathrm{HMT}},\varphi_{\mathrm{HMT}}\),
\(\delta\), and \(\lambda\), with
\begin{equation}
 0\leq F_j^+-F_j^-\leq \frac{c}{4}\,(s_j-r_j).
 \label{eq:alpha-anchura-evaluacion-rescate-rev08}
\end{equation}
When \(F_j^+<0\), the right half is retained; when \(F_j^->0\), the left
half is retained. If \(F_j^-\leq0\leq F_j^+\), no attempt is made to decide
whether the midpoint is exactly the root. Let \(w_j=s_j-r_j\). The uniform
bound \(E'_{\mathsf s,\infty}\geq c=6\) and the mean value theorem imply
that any root contained in \(D_j\) belongs to
\begin{equation}
 D_{j+1}=D_j\cap
 \left[m_j-\frac{w_j}{4},
             m_j+\frac{w_j}{4}\right].
 \label{eq:alpha-rescate-derivada-rev08}
\end{equation}
Indeed, let \(\xi_j\) be the root and let
\(y_j=E_{\mathsf s,\infty}(m_j)\in[F_j^-,F_j^+]\). Since the enclosure
contains zero and satisfies
\eqref{eq:alpha-anchura-evaluacion-rescate-rev08}, one has
\(|y_j|\leq cw_j/4\). The mean value theorem gives
\(|m_j-\xi_j|\leq |y_j|/c\leq w_j/4\), which proves
\eqref{eq:alpha-rescate-derivada-rev08}. The argument includes the case
\(y_j=0\): the root is then at the midpoint, but the algorithm need not
recognize that equality in order to retain it.

If the available enclosure does not satisfy
\eqref{eq:alpha-anchura-evaluacion-rescate-rev08}, the source cylinders are
first refined and the evaluation is repeated; no sign is declared and no
insufficient contraction is accepted. In each of the three branches, the new
interval retains the root and has diameter at most \(w_j/2\). For every
finite target depth, directed refinement makes it possible to satisfy the
width bound without making an uncertified comparison of a real number with
zero. Monotonicity proves inductively that
\begin{equation}
 a(\mathsf s)\in D_{j+1}\subseteq D_j,\qquad
 \operatorname{diam}D_{j+1}\leq
 \tfrac12\operatorname{diam}D_j,
 \qquad \operatorname{diam}D_j\longrightarrow0,
 \label{eq:alpha-biseccion-racional}
\end{equation}
where the inequality is a contraction bound, not an imposed equality. The
executable certificate applies this rule to twenty-four levels in base one
thousand and checks that the complete outer closure \(A_{1000}\), not merely
its lower endpoint, lies inside every expressed cylinder. Its negative tests
reject an enclosure that is too wide, a disordered enclosure, a point that is
not the midpoint, a domain with no sign change, and a degenerate bracket. The
coinductive prolongation permits one thousand to be replaced by any required
depth.

Concretely, for \(S_N=1000^N\), the verifier computes
\begin{equation}
 j_N^-:=\lfloor S_Na_-\rfloor,
 \qquad j_N^+:=\lfloor S_Na_+\rfloor
 \label{eq:alpha-indices-cilindro-dirigidos-rev08}
\end{equation}
and expresses the level only if \(j_N^-=j_N^+\). This equality, which
includes the upper endpoint of the outer closure, proves
\(A_{1000}\subset[j_N^-/S_N,(j_N^-+1)/S_N)\) and avoids assigning to the
left cell an endpoint lying exactly on its right boundary.

Let \(C_N^A(\mathsf s)\) be the integral cylinder in
\eqref{eq:alpha-cilindros-via-a-exactos}. Recursively choose
\(m(N)>m(N-1)\) so that \(\operatorname{diam}D_{m(N)}<1000^{-N}/4\), and
define
\begin{equation}
I_{B,N}:=D_{m(N)}\cap C_N^A(\mathsf s).
 \label{eq:alpha-intervalos-b-completos}
\end{equation}
Define \(\ell_N:=\inf I_{B,N}\) and \(u_N:=\sup I_{B,N}\). Both are
rational; the interval is nonempty because both factors contain
\(a(\mathsf s)\). Since the cylinder is half-open, \(u_N\) need not belong
to \(I_{B,N}\). Evaluation at that endpoint is understood through the
continuous extension of \(E_{\mathsf s,\infty}\) to the outer closure. Thus,
\begin{equation}
 I_{B,N+1}\subseteq I_{B,N}\subseteq C_N^A(\mathsf s),
 \qquad \operatorname{diam}I_{B,N}\longrightarrow0,
 \qquad E_{\mathsf s,\infty}(\ell_N)\leq0\leq
 E_{\mathsf s,\infty}(u_N).
 \label{eq:alpha-inclusion-signos-completa}
\end{equation}
Bound~\eqref{eq:alpha-cota-cola-uniforme} permits a truncation to be chosen
that preserves the sign at each endpoint where the complete function does not
vanish. If an endpoint coincides with the root, the identity of the complete
function and the membership already proved are used; that same sign is not
attributed to its truncations. Indeed, the exact remainder is
\[
 E_{\mathsf s,\infty}(x)-E_{\mathsf s,M}(x)
 =\lambda(\mathsf s)x^{10}
   \frac{(qx^3)^{M+1}}{1-qx^3}.
\]
In the positive domain under consideration, where \(\lambda(\mathsf s)<0\),
one obtains at the root
\[
 E_{\mathsf s,M}(a(\mathsf s))
 =-\lambda(\mathsf s)a(\mathsf s)^{10}
   \frac{(qa(\mathsf s)^3)^{M+1}}{1-qa(\mathsf s)^3}>0
 \qquad(M<\infty).
\]
Therefore, certification of the weak signs in
\eqref{eq:alpha-inclusion-signos-completa} refers to
\(E_{\mathsf s,\infty}\). Its evaluation uses the complete function or the
truncation accompanied by the signed remainder. This covers both the included
infimum and the supremum that may be excluded from the cylinder.

\begin{proposition}[Compatibility squares and common cylinders]
\label{prop:alpha-cuadrados-compatibilidad}
The restrictions of words and analytic representations commute:
\[
\begin{array}{ccc}
 \mathcal X_\alpha^{\mathrm{enr}}&
 \xrightarrow{\ \mathsf T_{\alpha,N'}\ }&\mathcal W_{N'}\\[1mm]
 \Big\Vert&&\Big\downarrow\rho_{N',N}\\[1mm]
 \mathcal X_\alpha^{\mathrm{enr}}&
 \xrightarrow{\ \mathsf T_{\alpha,N}\ }&\mathcal W_N
\end{array}
\qquad
\begin{array}{ccc}
 \mathcal X_\alpha^{\mathrm{enr}}&
 \xrightarrow{\ \Gamma_{M'}\ }&\mathfrak J_{M'}\\[1mm]
 \Big\Vert&&\Big\downarrow\tau_{M',M}\\[1mm]
 \mathcal X_\alpha^{\mathrm{enr}}&
 \xrightarrow{\ \Gamma_M\ }&\mathfrak J_M .
\end{array}
\]
The collection \(I_{B,N}\) provides the effective connection between the two
squares and materially satisfies \(I_{B,N}\subseteq C_N^A\) for every \(N\).
\end{proposition}

\begin{proof}
The first square is the compatibility of the integral normalization; the
second is Proposition~\ref{prop:alpha-naturalidad-recurrencia-explicita}.
The inclusion follows from Definition~\eqref{eq:alpha-intervalos-b-completos},
and nonemptiness and the signs follow from
Theorem~\ref{thm:alpha-raiz-completa-explicita}. No step identifies the root
of a truncation with that of an arbitrary tail.
\end{proof}

We now define, without leaving the prefix reader implicit,
\begin{equation}
 \alpha_B(\mathsf s):=
 \operatorname*{arg\,zero}_{x\in I_\alpha}E_{\mathsf s,\infty}(x),
 \qquad
 \operatorname{pref}_N(x):=
 \frac{\lfloor1000^Nx\rfloor}{1000^N},
 \qquad
 \alpha_{B,N}:=\operatorname{pref}_N(\alpha_B).
 \label{eq:alpha-definicion-b-y-prefijos-rev08}
\end{equation}

\begin{theorem}[Projective equality of the two correlated representations]
\label{prop:alpha-dos-limites}
For every \(N\), the integral and analytic routes express the same prefix:
\begingroup
\renewcommand{\theequation}{A7\(_0\)}
\begin{equation}
 \boxed{\alpha_{A,N}=\operatorname{pref}_N(\alpha_A)
 =\operatorname{pref}_N(\alpha_B)=\alpha_{B,N}.}
 \label{eq:alpha-dos-vias-cada-nivel}
\end{equation}
\endgroup
In the limit,
\begingroup
\renewcommand{\theequation}{A8\(_0\)}
\begin{equation}
 \boxed{\alpha_A
 =\operatorname*{arg\,zero}_{x\in I_\alpha}
       E_{\mathsf s,\infty}(x)
 =\alpha_B=\alpha_{\mathrm{HMT}}.}
 \label{eq:alpha-dos-vias-limite-coinductivo}
\end{equation}
\endgroup
\end{theorem}

\begin{proof}
The integral normalization expresses exactly \(a(\mathsf s)\).
Theorem~\ref{thm:alpha-raiz-completa-explicita} proves that the analytic chart
represents that same point as its unique simple root. The half-open cylinders
\(C_N^A\) fix a canonical prefix, and \(I_{B,N}\subseteq C_N^A\) contains
\(\alpha_B=\alpha_A\). Therefore, their prefixes agree for every \(N\); the
nesting and the diameters tending to zero give equality of the limits.
Metrological recognition occurs subsequently.
\end{proof}

The analytic function constitutes a correlated representation of the same
state, not a second derivation causally independent of the integral closure.
This precision does not weaken the proved equality: it identifies its exact
mechanism, prevents attribution to the jet of a tail that it does not
determine, and prevents promotion of a conventional value to a generator.
